\chapter{Conclusiones}

\section{Conclusiones principales}
\begin{itemize}
  \item El modelo con efecto Allee y no reciprocidad genera umbrales claros; en los barridos realizados todos los equilibrios con $i^*>0$ son inestables.
  \item El control variacional $\mu=1$ compacta patrones y reduce la heterogeneidad espacial, pero no elimina la inestabilidad de fondo.
  \item El efecto Allee fuerte limita la extensión tumoral e impone una barrera crítica; el efecto débil permite patrones más amplios y sensibles a control.
\end{itemize}

\section{Aportaciones originales}
\begin{itemize}
  \item Integración de efecto Allee (débil/fuerte), no reciprocidad y control inmunológico en un mismo marco ODE/PDE.
  \item Tablas de candidatos inestables y rejillas de parámetros para iniciar simulaciones PDE comparables.
  \item Análisis correlacional espacio--temporal que vincula control variacional con organización de patrones.
\end{itemize}

\section{Trabajo futuro}
\begin{itemize}
  \item Buscar equilibrios estables ajustando $rc,\beta,\eta$ y aumentando $rd,\delta$ con mallas más finas.
  \item Explorar protocolos adaptativos $u(x,y,t)$ y optimización de costes de control.
  \item Extender simulaciones a horizontes temporales mayores y geometrías no uniformes.
\end{itemize}




